\section{Results}

Unless stated otherwise, the reported contrasts are descriptive measurements of frozen native assignment states. Image-cluster intervals quantify sampling uncertainty conditional on the audited checkpoint. O2O--O2M contrasts compare states with different supervision semantics, first-divergence labels identify the earliest instrumented change, and cardinal distances follow four grid-resolved rays. These conventions govern the Results below and are repeated only where a local support or conditioning restriction changes the interpretation.

\subsection{Qualification of the native replay audit}

% Evidence: EV-INSTR-ORACLE-V1, EV-INSTR-MUTATION-V1, EV-INSTR-PARITY-5251-V1, EV-BOUNDARY-REFERENCE-FORMAL-V2
The qualification suite established that the replay instrument reproduces the native state before scientific interpretation. The analytical suite recovered all six oracle cases and seven property invariants, and it detected all six prespecified mutations. Native execution, production tracing, and the independent NumPy reference produced zero mismatch across 300 base states and 5,251 deterministic directional states. The independent implementation also reproduced all 4,408 continuous-boundary qualification trajectories under the then-current schedule. The focal-row engine then matched complete native replay at every state in 3,072 full-grid scenarios, with byte-identical event and curve files. Grid resolution was assessed separately and led to the full-domain $1/1024$ primary analysis.

\subsection{Stress parameterisation exposes branch-specific response structure}

% Evidence: EV-STRESS-CONTRACT-PRIMARY-V1
The apparent scale response changed substantially with the stress contract, while the two native branches remained descriptively distinct. On the AI-TOD-v2 seed-0 discovery sample, fixed one-pixel replay produced an O2O 8--16 minus 16--32 px fragility contrast of \FixedOtwoOScalePP\ pp. Under equivalent-area-side replay, the O2O contrast was \NormalizedOtwoOScalePP\ pp and its interval included zero. O2M instead changed from an uncertain fixed contrast of \FixedOtwoMScalePP\ pp to an equivalent-area-side contrast of \NormalizedOtwoMScalePP\ pp. The paired O2O-minus-O2M contrast was positive under both contracts (Table~\ref{tab:stress_primary}), but fixed stress was dominated by the O2O response whereas equivalent-area-side stress was dominated by the reverse O2M response. The two contracts were not pooled because their displacements answer different questions.

The high absolute O2M assigned-set Top-1 fragility rates in Table~\ref{tab:stress_primary} are object-level fragility estimates, not directional flip rates, and should not be interpreted as loss of positive supervision. Assigned-set Top-1 is a rank probe inside a multi-positive state; it is not equivalent to loss of positive supervision. Figure~\ref{fig:rank_set_geometry} separately measures whether the positive set itself persisted.

Direction-normalised replay changed the O2O gap to 2.79 pp [0.44, 4.96], while the O2M gap remained $-17.04$ pp [$-22.16$, $-11.58$] and the paired gap was 19.83 pp [13.65, 25.83]. A separate paired composition sensitivity standardised both size groups to the same aspect-ratio distribution; its paired gaps were 17.53 pp [11.07, 24.08] for equivalent-area-side replay and 19.72 pp [13.67, 25.77] for direction-normalised replay. Rectangular-box parameterisation therefore affected the absolute O2O response, with the largest point difference among elongated objects, but did not remove the paired branch pattern.

\begin{table*}[t]
\centering
\caption{Primary AI-TOD-v2 fragility rates and stress-contract contrasts in percentage points. Brackets are 95\% stratified image-cluster bootstrap intervals. Gaps are 8--16 minus 16--32 px; positive values denote higher fragility in the 8--16 px group. The paired gap is O2O minus O2M. These are discovery estimates. Support: fixed primary, 6,378 common-valid GTs from 288 contributing images; equivalent-area primary, 6,377 common-valid GTs from 288 contributing images.}
\label{tab:stress_primary}
\begin{tabular}{llrrrr}
\toprule
Stress contract & Branch & 8--16 rate & 16--32 rate & Scale gap [95\% CI] & Paired gap [95\% CI] \\
\midrule
Fixed 1 px & O2O & 22.14 & 7.77 & 14.37 [10.84, 17.78] & 12.95 [6.52, 19.16] \\
Fixed 1 px & O2M & 71.74 & 70.32 & 1.42 [-3.18, 6.31] & -- \\
Equivalent-area side $\kappa=0.0625$ & O2O & 9.11 & 8.98 & 0.13 [-3.41, 3.15] & 17.58 [11.06, 23.98] \\
Equivalent-area side $\kappa=0.0625$ & O2M & 57.93 & 75.37 & -17.45 [-22.72, -11.71] & -- \\
\bottomrule
\end{tabular}

\end{table*}

% Evidence: EV-STRESS-CONTRACT-KAPPA-DATASET-V1
Across the three-point equivalent-area-side grid, all six O2M contrasts were negative and all six paired branch contrasts were positive. On AI-TOD-v2, paired contrasts ranged from 17.27 to 18.19 pp; on VisDrone they ranged from 34.79 to 37.25 pp. Absolute O2O contrasts varied with both $\kappa$ and dataset: the AI-TOD-v2 intervals at $\kappa=0.03125$ and $0.0625$ included zero, whereas the VisDrone contrasts were positive at all three magnitudes. The paired direction was stable over the tested grid, whereas the absolute O2O response was not.

% Evidence: EV-YOLOV10-NORMALIZED-V1
The compatible YOLOv10-S contract showed the same equivalent-area-side branch pattern: O2O was \YtenOtwoOPP\ pp with an interval spanning zero, O2M was \YtenOtwoMPP\ pp, and the paired contrast was \YtenPairedPP\ pp. This provides replication across two compatible detector contracts without establishing architecture invariance.

% Evidence: EV-LEARNING-DYNAMICS-V1
The branch pattern was present at all six physically archived stages spanning the first completed epoch to the final training state. Across those stages, $\mathrm{mAP}_{50:95}$ changed from \DynamicsMapLow\ to \DynamicsMapHigh, while fixed paired contrasts ranged from \DynamicsFixedPairedLowPP\ to \DynamicsFixedPairedHighPP\ pp and equivalent-area-side paired contrasts from \DynamicsNormPairedLowPP\ to \DynamicsNormPairedHighPP\ pp. No checkpoint was available between epochs 100 and 280, so no continuity or interpolation is claimed across that interval. This six-stage archive is distinct from the retained files named \texttt{best.pt}: its final row uses the archived \texttt{last.pt} training state labelled 300, whereas the separate seed sensitivity uses post-training-selected artifacts with embedded zero-based epoch 280.

% Evidence: EV-MULTISEED-FIXED-COMPLETION-V1, EV-MULTISEED-EQUIVALENT-RANKSET-V1
Across every protocol-eligible released seed, the fixed paired contrast on the strict 6,318-GT intersection was 12.47, 12.44, 13.19, and 12.88 pp for seeds 0--3; the equal-seed mean was 12.75 pp [6.27, 18.84]. On the strict 6,317-GT equivalent-area intersection, paired contrasts were 17.20, 17.51, 17.87, and 19.18 pp; the equal-seed mean was 17.94 pp [11.51, 24.23]. Equivalent-area O2M gaps were negative in all four seeds, and the positive-set-turnover contrast remained positive in all four. Both 20,000-replicate crossed seed and image-cluster bootstraps therefore retained the branch direction at the two endpoints and the Rank--Set direction across the eligible checkpoints. These are $S=4$ sensitivities; eligibility, pathway, and boundary-distance estimates remain conditional on seed 0.

Figure~\ref{fig:stress_response} summarises these four sensitivity axes. Panel C uses categorical positions and marks the absent epoch 100--280 archives; no interpolation is implied.

\begin{figure*}[t]
\centering
\includegraphics[width=\textwidth]{final_v4/Figure_2.pdf}
\caption{The O2O-minus-O2M contrast remained positive across the audited stress magnitudes, datasets, archived checkpoints, and detector contracts, while branch-wise responses differed between stress contracts. (A) AI-TOD-v2 O2O, O2M, and paired scale contrasts under fixed one-pixel and equivalent-area-side replay. (B) Paired contrasts at the three prespecified equivalent-area-side $\kappa$ values on AI-TOD-v2 and VisDrone. The values occupy their true positions on a linear numerical axis; no connecting curve or continuous trend model is fitted. (C) Paired contrasts at six equally spaced categorical archived checkpoints labelled 0, 20, 40, 60, 80, and 300 (final); no checkpoints were archived for epochs 100--280. (D) Equivalent-area-side replication under compatible YOLO26s and YOLOv10-S detector contracts. Points are inverse-inclusion-weighted estimates and whiskers are 95\% stratified image-cluster bootstrap intervals.}
\label{fig:stress_response}
\end{figure*}

\subsection{Eligibility locking and first-observed pathway localisation}

% Evidence: EV-ELIGIBILITY-LOCK-V1
Hard eligibility boundaries accounted for much of the fixed-pixel response on the lock-analysis support, while equivalent-area-side replay shifted pathway localisation toward within-set rank changes. The lock analysis used 7,434 fixed-arm objects from 300 selected images, whereas the primary estimate used 6,378 common-valid O2O/O2M objects from 288 contributing images. On the lock support, the native fixed-pixel O2O gap was 17.10 pp and the locked gap was 3.52 pp. Their difference was \EligibilityLockFixedGapPP\ pp, corresponding to \EligibilityLockFixedFractionPct\% [70.51\%, 91.92\%] of the 17.10 pp denominator, not of the separate 14.37 pp primary gap. The equivalent-area-side reduction was \EligibilityLockNormalizedGapPP\ pp [1.42, 4.83]. The 2.95-pp native equivalent-area value is a property of the independent lock-analysis support and is not a re-estimate of the 0.13-pp common-valid primary contrast.

% Evidence: EV-ELIGIBILITY-STABLE-V2
The observational eligibility-stable subset retained 6,217 of 7,434 fixed-replay objects and 6,999 of 7,432 equivalent-area-side objects. Its fixed O2O scale contrast was 3.36 pp [1.11, 5.48], whereas its equivalent-area-side contrast was 1.05 pp [$-1.19$, 2.96]. The residual fixed response was smaller than the native response but remained positive. Because eligibility stability is conditioned after replay, these estimates have a different, potentially selected estimand.

% Evidence: EV-PATHWAY-ANATOMY-V2
First-divergence anatomy showed the corresponding change in pathway composition. Eligibility was the first observed difference for \FixedEligibilityPct\% [63.02\%, 71.07\%] of fixed-pixel identity exchanges and \NormalizedEligibilityPct\% [46.06\%, 55.94\%] under equivalent-area-side replay. Within-set geometric rank reversal was first for \FixedWithinPct\% [26.14\%, 34.36\%] and \NormalizedWithinPct\% [40.62\%, 50.26\%], respectively. Because first divergence records the earliest stage in a fixed hierarchy, these shares are order-dependent localisations.

The nonexclusive audit showed substantial downstream co-occurrence. Under fixed replay, a within-set rank change occurred on 75.12\% [70.98\%, 79.35\%] of identity-exchange paths. It was first on only 30.13\%, and eligibility co-occurred on 44.63\% [41.07\%, 48.81\%]. Under equivalent-area-side replay, any-stage within-set change was 77.44\% [73.15\%, 81.84\%]. It was first on 45.37\%, and eligibility co-occurred on 31.59\% [27.78\%, 35.68\%]. The gap between first-label and any-stage occurrence shows how pipeline order changes the localisation summary.

% Evidence: EV-PRERESOLUTION-FIXED-V1, EV-PRERESOLUTION-NORMALIZED-V1
The aggregate O2O branch response was already present before final conflict resolution. The assigned-minus-pre-conflict scale-gap difference was $-0.077$ pp [$-0.367$, 0.313] for fixed replay. Under equivalent-area-side replay, it was $-0.395$ pp [$-0.986$, 0.034]. Object-level fragility disagreement was 0.578\% under both contracts. All 134 fixed directional disagreements and all 96 equivalent-area-side disagreements coincided with the instrumented conflict stage. Conflict localised this small set of disagreements but did not explain the aggregate branch-scale pattern.

These boundary-localisation results are shown in Figure~\ref{fig:pathway_localisation}. The post-replay eligibility-stable subset remains a distinct selected estimand and is reported numerically in the Supplementary Material rather than repeated in the main figure.

\begin{figure*}[t]
\centering
\includegraphics[width=\textwidth]{final_v4/Figure_3.pdf}
\caption{Eligibility locking removed most of the fixed-pixel O2O contrast, while within-set changes occurred more often anywhere along an exchange path than exclusive first labels suggested. (A) Native and eligibility-locked O2O contrasts on the same lock-analysis support. (B) Exclusive first-divergence composition (left) and nonexclusive any-stage occurrence for eligibility, Top-$k$, conflict, within-set rank, and eligibility--rank co-occurrence (right). (C) Pre-conflict Top-7 winner and final post-conflict assigned contrasts provide a same-branch pre-resolution control. Filled markers denote fixed one-pixel replay and open markers equivalent-area-side replay. Whiskers are 95\% stratified image-cluster bootstrap intervals. First divergence localises the earliest observed difference and is not a causal contribution.}
\label{fig:pathway_localisation}
\end{figure*}

\subsection{O2O and O2M occupy different cardinal boundary structures}

% Evidence: EV-BOUNDARY-GRID-CONVERGENCE-V1, EV-BOUNDARY-GEOMETRY-V3
Assigned O2O identity and O2M assigned-set Top-1 showed different first-departure persistence along the audited cardinal rays. The analysis included 29,724 directional trajectories from 7,431 focal GTs. Through $\tau=0.125$, O2O had an RMCBD of 0.11985 (95\% CI 0.11932--0.12032), versus 0.10025 (0.09864--0.10180) for O2M. The O2O-minus-O2M contrast was 0.01745 (0.01586--0.01898) for 8--16 px objects and 0.03319 (0.02937--0.03652) for 16--32 px objects. By $\tau$, O2O boundaries had been observed for 9.05\% (8.28--9.92\%) of tiny and 6.39\% (4.61--8.49\%) of small trajectories; the corresponding O2M rates were 32.75\% (31.27--34.22\%) and 48.25\% (44.89--51.40\%). The branches differed in both observed cardinal persistence and scale structure. Exact RMCBD values and intervals are tabulated in the Supplementary Material.

Figure~\ref{fig:boundary_geometry} summarises full-domain $1/1024$ identity survival, cause-specific incidence, restricted mean distance, and truncation-horizon sensitivity.

\begin{figure*}[t]
\centering
\includegraphics[width=\textwidth]{final_v4/Figure_4.pdf}
\caption{Post-conflict assigned O2O identities persisted farther along cardinal rays than O2M assigned-positive Top-1 identities across the full-domain $1/1024$ analysis. (A) Weighted Kaplan--Meier summaries of remaining free of a first observed state departure under the recorded administrative and geometric censoring rules, shown as first-departure survival $S(r)$ for O2O and O2M in the 8--16 and 16--32 px groups. (B) Aalen--Johansen cumulative incidence of the first observed O2O boundary causes. (C) Absolute branch-specific restricted mean cardinal boundary distance (RMCBD) and paired O2O-minus-O2M contrasts through the primary horizon $\tau=0.125$. (D) Paired contrasts across six prespecified horizons on a linear numerical $\tau$ axis. Each branch-specific RMCBD is nondecreasing in $\tau$ by construction, whereas monotonicity of their difference is not guaranteed; the paired contrast increased over the observed range. Distances follow four cardinal rays and are not two-dimensional minimum boundaries. KM/AJ/RMCBD are descriptive summaries under the recorded censoring rules.}
\label{fig:boundary_geometry}
\end{figure*}

\subsection{O2M Rank--Set turnover and positive-count conditioning}

% Evidence: EV-RANK-SET-FIXED-V1, EV-RANK-SET-NORMALIZED-V1
O2M Top-1 turnover and positive-set turnover were distinct native states under both stress contracts. In the canonical 16--32 minus 8--16 px orientation, fixed replay gave Jaccard and base-retention contrasts of $+4.19$ and $+2.26$ pp, whereas equivalent-area-side replay gave $-5.91$ and $-4.93$ pp. The equivalent-area-side assigned-set Top-1 flip contrast was $+12.80$ pp, while the rank-only share among flips was $-22.50$ pp. Smaller objects had fewer Top-1 flips but a larger rank-only fraction among those flips; larger objects combined greater Top-1 turnover with lower set persistence.

% Evidence: EV-POSITIVE-COUNT-THREE-STAGE-V1
The three-stage count analysis separated support restriction from count-composition standardisation. Native positive count averaged 3.07 for 8--16 px and 6.54 for 16--32 px objects. Restricting the object-level Top-1 fragility analysis from full support to the adequate exact-$K=3$--7 overlap changed the 16--32-minus-8--16 gap from 17.45 pp [11.85, 22.73] to 4.53 pp [$-1.46$, 10.52]; standardising $K$ composition within that overlap changed it further to 2.16 pp [$-4.74$, 9.71]. The direction-level sequence was 11.59 pp [8.42, 14.66], 3.59 pp [0.80, 6.63], and 0.88 pp [$-2.18$, 4.47]. The overlap covered 60.3\% and 64.3\% of the two size groups. Positive-set-turnover gaps instead increased from 33.16 pp [31.35, 34.79] on full support to 36.82 pp [34.73, 38.87] on overlap support and 41.13 pp [38.23, 44.26] after standardisation. Most Top-1 attenuation therefore arose from the support restriction; the additional within-overlap standardisation effect was smaller. Neither change identifies a causal cardinality effect, and the frozen tables lack the complete $q$ ordering required for a genuine fixed-$K$ assigner replay.

Figure~\ref{fig:rank_set_geometry} separates rank turnover from set persistence and shows how count conditioning affects their scale contrasts.

\begin{figure*}[t]
\centering
\includegraphics[width=\textwidth]{final_v4/Figure_5.pdf}
\caption{Support restriction accounted for most attenuation of the O2M Top-1 contrast, while positive-set turnover remained large. (A) Schematic distinction among rank-only turnover, simultaneous rank-and-set turnover, and set-only turnover with stable assigned-set Top-1. Filled circles denote Top-1 and open circles other positives. (B) Native O2M Rank--Set contrasts under fixed and equivalent-area-side replay, grouped into set persistence and rank turnover. (C) Weighted native positive-count distributions; $K=0$ remains visible and shading marks the $K=3$--7 overlap used for standardisation. (D) Full-support crude, overlap-support crude, and overlap-support count-standardised contrasts. Panels B and D use 16--32 minus 8--16 px throughout. The three markers separate support restriction from within-overlap $K$ standardisation; neither operation identifies a causal count intervention.}
\label{fig:rank_set_geometry}
\end{figure*}

% Evidence: EV-SAME-STRIDE-V1
The strict same-stride analysis retained 3,835 paired GTs. Its O2O scale contrast was $+10.23$ pp [7.30, 13.32], its O2M contrast was $-0.12$ pp [$-6.83$, 5.78], and the paired difference was $+10.35$ pp [3.54, 18.17]. Cross-stride switching alone therefore did not account for the branch contrast, although post-replay conditioning can induce selection and does not exclude a feature-pyramid-level contribution.

\subsection{Native margin aligns with observed rank separation, with a narrower utility boundary}

% Evidence: EV-MARGIN-RHOR-CONSTRUCT-V3
Assigned-relative margin aligned with observed fixed-pair rank distance, but this conditional association did not translate into stable incremental diagnostic value. Only 1,593 of 29,724 trajectories had an observed fixed-pair $q$-order crossing; 7,439 were competing-censored, 16,548 right-censored, and 4,144 undefined. Within the observed-crossing subset, margin correlated with crossing distance at $\rho=0.628$ [0.584, 0.669]. Its correlation with observed eligibility-boundary distance was $\rho=0.150$ [0.111, 0.191], a difference of 0.477 [0.408, 0.543]. These are conditional observed-case associations rather than censoring-adjusted population correlations.

% Evidence: EV-MARGIN-SHAPE-SENSITIVITY-V1
Replacing a linear margin term with a training-fold-fitted restricted cubic spline did not improve held-out metrics consistently. For fragility, spline-minus-linear ROC-AUC and PR-AUC differences were $-0.00756$ [$-0.01411$, $-0.00138$] and $-0.00415$ [$-0.00813$, $-0.00109$]. Coverage-based-miss PR-AUC was lower by $-0.00541$ [$-0.01298$, $-0.00027$]. The log-loss intervals spanned zero for both outcomes. This negative sensitivity limits claims that a more flexible margin shape supplies additional diagnostic utility.

Figure~\ref{fig:margin_boundary} separates margin's observed-crossing rank-separation alignment from its more limited incremental diagnostic utility.

\begin{figure*}[t]
\centering
\includegraphics[width=\textwidth]{final_v4/Figure_6.pdf}
\caption{Assigned-relative margin aligned with observed fixed-pair rank-crossing distance but supplied no stable held-out diagnostic improvement. (A) Assigned-relative margin versus observed rank-crossing distance among 1,593 trajectories with an observed crossing. The line joins ten margin-decile medians as a descriptive summary; no regression curve is fitted. (B) Weighted composition of observed crossing, competing censoring, right censoring, and undefined states. (C) Conditional Spearman associations with observed rank crossing and eligibility-boundary distance on their separate observed supports. (D) Held-out diagnostic changes after adding margin on AI-TOD-v2 and VisDrone; log-loss and Brier signs are oriented so rightward values indicate improvement. These analyses concern conditional association and diagnostic increment, not a guaranteed predictor or intervention effect.}
\label{fig:margin_boundary}
\end{figure*}

\subsection{Locality and practical audit scope}

% Evidence: EV-SPILLOVER-FORMAL300-V1
The remaining analyses delimit the audit's locality, screening budget, and prospective evidence status. Focal replay changed a non-focal O2O identity in 52 of \SpilloverPairs\ pair-events. The overall weighted pair rate was 0.001465\% [0.000824\%, 0.002383\%]. All observed changes occurred among the 5,679 pairs sharing at least one native pre-conflict candidate claim; their rate was \SpilloverSharedPct\% [0.5787\%, 1.3248\%]. No nonsharing event was observed in this sample.

% Evidence: EV-SAMPLE-EFFICIENCY-1000-V1
At a 100-image audit budget, all 1,000 subsamples recovered the sign and O2O-versus-O2M branch ordering of the full equivalent-area-side paired contrast. The median absolute error was \SampleHundredMedianErrorPP\ pp and the 95th-percentile error was \SampleHundredPNinetyFiveErrorPP\ pp. Reduced budgets can screen a dominant response signature, whereas precise effect estimation still requires the full design.
